% Options for packages loaded elsewhere
\PassOptionsToPackage{unicode}{hyperref}
\PassOptionsToPackage{hyphens}{url}
\PassOptionsToPackage{dvipsnames,svgnames,x11names}{xcolor}
\documentclass[
  11pt,
  a4paper,
]{article}
\usepackage{xcolor}
\usepackage[margin=1in]{geometry}
\usepackage{amsmath,amssymb}
\setcounter{secnumdepth}{-\maxdimen} % remove section numbering
\usepackage{iftex}
\ifPDFTeX
  \usepackage[T1]{fontenc}
  \usepackage[utf8]{inputenc}
  \usepackage{textcomp} % provide euro and other symbols
\else % if luatex or xetex
  \usepackage{unicode-math} % this also loads fontspec
  \defaultfontfeatures{Scale=MatchLowercase}
  \defaultfontfeatures[\rmfamily]{Ligatures=TeX,Scale=1}
\fi
\usepackage{lmodern}
\ifPDFTeX\else
  % xetex/luatex font selection
  \ifLuaTeX
    \usepackage{luaotfload}
    \directlua{luaotfload.add_fallback("mainfontfallback",{
      "NotoSerifBengali:mode=harf"
    })}
  \fi
  \setmainfont[,RawFeature={fallback=mainfontfallback}]{Latin Modern Roman}
  \ifLuaTeX
    \usepackage{luaotfload}
    \directlua{luaotfload.add_fallback("sansfontfallback",{
      "NotoSansBengali:mode=harf"
    })}
  \fi
  \setsansfont[,RawFeature={fallback=sansfontfallback}]{Latin Modern Sans}
  \ifLuaTeX
    \usepackage{luaotfload}
    \directlua{luaotfload.add_fallback("monofontfallback",{
      "NotoSansBengali:mode=harf"
    })}
  \fi
  \setmonofont[,RawFeature={fallback=monofontfallback}]{DejaVu Sans Mono}
\fi
% Use upquote if available, for straight quotes in verbatim environments
\IfFileExists{upquote.sty}{\usepackage{upquote}}{}
\IfFileExists{microtype.sty}{% use microtype if available
  \usepackage[]{microtype}
  \UseMicrotypeSet[protrusion]{basicmath} % disable protrusion for tt fonts
}{}
\makeatletter
\@ifundefined{KOMAClassName}{% if non-KOMA class
  \IfFileExists{parskip.sty}{%
    \usepackage{parskip}
  }{% else
    \setlength{\parindent}{0pt}
    \setlength{\parskip}{6pt plus 2pt minus 1pt}}
}{% if KOMA class
  \KOMAoptions{parskip=half}}
\makeatother
\usepackage{color}
\usepackage{fancyvrb}
\newcommand{\VerbBar}{|}
\newcommand{\VERB}{\Verb[commandchars=\\\{\}]}
\DefineVerbatimEnvironment{Highlighting}{Verbatim}{commandchars=\\\{\}}
% Add ',fontsize=\small' for more characters per line
\usepackage{framed}
\definecolor{shadecolor}{RGB}{248,248,248}
\newenvironment{Shaded}{\begin{snugshade}}{\end{snugshade}}
\newcommand{\AlertTok}[1]{\textcolor[rgb]{0.94,0.16,0.16}{#1}}
\newcommand{\AnnotationTok}[1]{\textcolor[rgb]{0.56,0.35,0.01}{\textbf{\textit{#1}}}}
\newcommand{\AttributeTok}[1]{\textcolor[rgb]{0.13,0.29,0.53}{#1}}
\newcommand{\BaseNTok}[1]{\textcolor[rgb]{0.00,0.00,0.81}{#1}}
\newcommand{\BuiltInTok}[1]{#1}
\newcommand{\CharTok}[1]{\textcolor[rgb]{0.31,0.60,0.02}{#1}}
\newcommand{\CommentTok}[1]{\textcolor[rgb]{0.56,0.35,0.01}{\textit{#1}}}
\newcommand{\CommentVarTok}[1]{\textcolor[rgb]{0.56,0.35,0.01}{\textbf{\textit{#1}}}}
\newcommand{\ConstantTok}[1]{\textcolor[rgb]{0.56,0.35,0.01}{#1}}
\newcommand{\ControlFlowTok}[1]{\textcolor[rgb]{0.13,0.29,0.53}{\textbf{#1}}}
\newcommand{\DataTypeTok}[1]{\textcolor[rgb]{0.13,0.29,0.53}{#1}}
\newcommand{\DecValTok}[1]{\textcolor[rgb]{0.00,0.00,0.81}{#1}}
\newcommand{\DocumentationTok}[1]{\textcolor[rgb]{0.56,0.35,0.01}{\textbf{\textit{#1}}}}
\newcommand{\ErrorTok}[1]{\textcolor[rgb]{0.64,0.00,0.00}{\textbf{#1}}}
\newcommand{\ExtensionTok}[1]{#1}
\newcommand{\FloatTok}[1]{\textcolor[rgb]{0.00,0.00,0.81}{#1}}
\newcommand{\FunctionTok}[1]{\textcolor[rgb]{0.13,0.29,0.53}{\textbf{#1}}}
\newcommand{\ImportTok}[1]{#1}
\newcommand{\InformationTok}[1]{\textcolor[rgb]{0.56,0.35,0.01}{\textbf{\textit{#1}}}}
\newcommand{\KeywordTok}[1]{\textcolor[rgb]{0.13,0.29,0.53}{\textbf{#1}}}
\newcommand{\NormalTok}[1]{#1}
\newcommand{\OperatorTok}[1]{\textcolor[rgb]{0.81,0.36,0.00}{\textbf{#1}}}
\newcommand{\OtherTok}[1]{\textcolor[rgb]{0.56,0.35,0.01}{#1}}
\newcommand{\PreprocessorTok}[1]{\textcolor[rgb]{0.56,0.35,0.01}{\textit{#1}}}
\newcommand{\RegionMarkerTok}[1]{#1}
\newcommand{\SpecialCharTok}[1]{\textcolor[rgb]{0.81,0.36,0.00}{\textbf{#1}}}
\newcommand{\SpecialStringTok}[1]{\textcolor[rgb]{0.31,0.60,0.02}{#1}}
\newcommand{\StringTok}[1]{\textcolor[rgb]{0.31,0.60,0.02}{#1}}
\newcommand{\VariableTok}[1]{\textcolor[rgb]{0.00,0.00,0.00}{#1}}
\newcommand{\VerbatimStringTok}[1]{\textcolor[rgb]{0.31,0.60,0.02}{#1}}
\newcommand{\WarningTok}[1]{\textcolor[rgb]{0.56,0.35,0.01}{\textbf{\textit{#1}}}}
\usepackage{longtable,booktabs,array}
\usepackage{calc} % for calculating minipage widths
% Correct order of tables after \paragraph or \subparagraph
\usepackage{etoolbox}
\makeatletter
\patchcmd\longtable{\par}{\if@noskipsec\mbox{}\fi\par}{}{}
\makeatother
% Allow footnotes in longtable head/foot
\IfFileExists{footnotehyper.sty}{\usepackage{footnotehyper}}{\usepackage{footnote}}
\makesavenoteenv{longtable}
\setlength{\emergencystretch}{3em} % prevent overfull lines
\providecommand{\tightlist}{%
  \setlength{\itemsep}{0pt}\setlength{\parskip}{0pt}}
% paper_preamble.tex — shared LaTeX preamble for the PCMA/SI paper set.
%
% Compiled with lualatex (NOT pdflatex) because several papers contain
% Bengali script: the CLL paper realises one cognitive state into
% English, Bengali and Banglish, so the Bengali output is EVIDENCE and
% cannot be transliterated away without destroying the point being made.
%
% Font selection is NOT done here. pandoc >= 3 has native
% `mainfontfallback` / `monofontfallback` variables that emit the
% luaotfload fallback registration in the correct position relative to
% unicode-math; doing it by hand in this header lands after unicode-math
% and breaks font loading outright. See tools/build_papers.sh.

\usepackage{amsmath}
\usepackage{amssymb}
\usepackage{mathtools}   % \xleftrightarrow etc. used in cll.md §5
\usepackage{booktabs}

% ---------------------------------------------------------------------
% Unicode symbols appearing in PROSE rather than inside $...$.
%
% The papers were written as markdown and read as markdown, so they use
% real characters — "the ω-boost", "α-child", "ΔR coherence" — instead of
% LaTeX macros. Under unicode-math those codepoints are math-mode-only
% and abort the build ("\symbballowed only in math mode"). Mapping each
% to \ensuremath keeps the markdown readable on GitHub AND compilable,
% which is the whole point of keeping one source.
% ---------------------------------------------------------------------
\usepackage{newunicodechar}

% Greek, lower case
\newunicodechar{α}{\ensuremath{\alpha}}
\newunicodechar{β}{\ensuremath{\beta}}
\newunicodechar{γ}{\ensuremath{\gamma}}
\newunicodechar{δ}{\ensuremath{\delta}}
\newunicodechar{ε}{\ensuremath{\varepsilon}}
\newunicodechar{η}{\ensuremath{\eta}}
\newunicodechar{θ}{\ensuremath{\theta}}
\newunicodechar{κ}{\ensuremath{\kappa}}
\newunicodechar{λ}{\ensuremath{\lambda}}
\newunicodechar{μ}{\ensuremath{\mu}}
\newunicodechar{π}{\ensuremath{\pi}}
\newunicodechar{ρ}{\ensuremath{\rho}}
\newunicodechar{σ}{\ensuremath{\sigma}}
\newunicodechar{τ}{\ensuremath{\tau}}
\newunicodechar{φ}{\ensuremath{\varphi}}
\newunicodechar{χ}{\ensuremath{\chi}}
\newunicodechar{ψ}{\ensuremath{\psi}}
\newunicodechar{ω}{\ensuremath{\omega}}
\newunicodechar{ν}{\ensuremath{\nu}}
\newunicodechar{ξ}{\ensuremath{\xi}}
\newunicodechar{ζ}{\ensuremath{\zeta}}
\newunicodechar{ι}{\ensuremath{\iota}}

% Greek, upper case
\newunicodechar{Δ}{\ensuremath{\Delta}}
\newunicodechar{Φ}{\ensuremath{\Phi}}
\newunicodechar{Σ}{\ensuremath{\Sigma}}
\newunicodechar{Ψ}{\ensuremath{\Psi}}
\newunicodechar{Ω}{\ensuremath{\Omega}}
\newunicodechar{Γ}{\ensuremath{\Gamma}}
\newunicodechar{Λ}{\ensuremath{\Lambda}}
\newunicodechar{Θ}{\ensuremath{\Theta}}
\newunicodechar{Π}{\ensuremath{\Pi}}

% Operators and relations
\newunicodechar{×}{\ensuremath{\times}}
\newunicodechar{·}{\ensuremath{\cdot}}
\newunicodechar{±}{\ensuremath{\pm}}
\newunicodechar{−}{\ensuremath{-}}
\newunicodechar{≈}{\ensuremath{\approx}}
\newunicodechar{≤}{\ensuremath{\leq}}
\newunicodechar{≥}{\ensuremath{\geq}}
\newunicodechar{≠}{\ensuremath{\neq}}
\newunicodechar{∈}{\ensuremath{\in}}
\newunicodechar{∂}{\ensuremath{\partial}}
\newunicodechar{∇}{\ensuremath{\nabla}}
\newunicodechar{√}{\ensuremath{\sqrt{}}}
\newunicodechar{∞}{\ensuremath{\infty}}
\newunicodechar{∝}{\ensuremath{\propto}}
\newunicodechar{∑}{\ensuremath{\sum}}
\newunicodechar{∫}{\ensuremath{\int}}
\newunicodechar{ℝ}{\ensuremath{\mathbb{R}}}
\newunicodechar{ℕ}{\ensuremath{\mathbb{N}}}
\newunicodechar{‖}{\ensuremath{\|}}
\newunicodechar{⊕}{\ensuremath{\oplus}}
\newunicodechar{⊗}{\ensuremath{\otimes}}
\newunicodechar{∀}{\ensuremath{\forall}}
\newunicodechar{∃}{\ensuremath{\exists}}
\newunicodechar{¬}{\ensuremath{\neg}}
\newunicodechar{∧}{\ensuremath{\wedge}}
\newunicodechar{∨}{\ensuremath{\vee}}

% Arrows
\newunicodechar{→}{\ensuremath{\rightarrow}}
\newunicodechar{←}{\ensuremath{\leftarrow}}
\newunicodechar{↔}{\ensuremath{\leftrightarrow}}
\newunicodechar{↑}{\ensuremath{\uparrow}}
\newunicodechar{↓}{\ensuremath{\downarrow}}
\newunicodechar{⇒}{\ensuremath{\Rightarrow}}
\newunicodechar{⇐}{\ensuremath{\Leftarrow}}
\newunicodechar{⇔}{\ensuremath{\Leftrightarrow}}

% Super/subscripts written as literal characters
\newunicodechar{²}{\ensuremath{^{2}}}
\newunicodechar{³}{\ensuremath{^{3}}}
\newunicodechar{⁴}{\ensuremath{^{4}}}
\newunicodechar{⁵}{\ensuremath{^{5}}}
\newunicodechar{⁶}{\ensuremath{^{6}}}
\newunicodechar{₀}{\ensuremath{_{0}}}
\newunicodechar{₁}{\ensuremath{_{1}}}
\newunicodechar{₂}{\ensuremath{_{2}}}
\newunicodechar{₃}{\ensuremath{_{3}}}
\newunicodechar{ᵀ}{\ensuremath{^{\mathsf{T}}}}

% Status glyphs. The papers use emoji as status markers in tables
% (implemented / failed / warning). Mapped to typographic equivalents:
% a colour emoji font in lualatex is unreliable, and the MEANING is what
% carries — a reader needs to see "this one failed", not a red cross.
\newunicodechar{✅}{\ensuremath{\checkmark}}
\newunicodechar{✓}{\ensuremath{\checkmark}}
\newunicodechar{❌}{\ensuremath{\times}}
\newunicodechar{✗}{\ensuremath{\times}}
\newunicodechar{⚠}{\textbf{!}}
\newunicodechar{🧩}{\textbf{[piece]}}
\newunicodechar{📋}{\textbf{[list]}}
\newunicodechar{🟡}{\textbf{[partial]}}
\newunicodechar{🔴}{\textbf{[dead]}}
\newunicodechar{🟢}{\textbf{[live]}}
\newunicodechar{▼}{\ensuremath{\blacktriangledown}}
\newunicodechar{►}{\ensuremath{\blacktriangleright}}
\newunicodechar{□}{\ensuremath{\square}}
\newunicodechar{█}{\rule{0.6em}{0.9ex}}
\newunicodechar{•}{\textbullet}
\newunicodechar{†}{\textdagger}
\newunicodechar{‑}{-}
\newunicodechar{‐}{-}
\newunicodechar{…}{\ldots}
\newunicodechar{।}{\textbar}   % danda: rendered by the Bengali fallback in text
\newunicodechar{⁻}{\ensuremath{^{-}}}

% Zero-width formatting codepoints. These are invisible in an editor but
% are real characters to LaTeX: an emoji written as "⚠️" is TWO
% codepoints (U+26A0 U+FE0F), so declaring only the visible one leaves a
% variation selector with no definition and the build dies on a symbol
% nobody can see. Both are mapped to nothing.
\newunicodechar{^^^^fe0f}{}    % VARIATION SELECTOR-16
\newunicodechar{^^^^200d}{}    % ZERO WIDTH JOINER
\newunicodechar{^^^^200b}{}    % ZERO WIDTH SPACE

\usepackage{longtable}
\usepackage{array}

% Comparison matrices and cost ledgers are wide. Let them shrink to the
% text block rather than overflow the margin.
\usepackage{adjustbox}

% Code blocks: wrap long lines instead of running off the page. Several
% papers quote real C++ with long identifiers.
\usepackage{fvextra}
\fvset{breaklines=true,breakanywhere=true,fontsize=\small}

\setlength{\parskip}{0.5em}
\setlength{\parindent}{0pt}

\usepackage{fancyhdr}
\pagestyle{fancy}
\fancyhf{}
\fancyfoot[C]{\thepage}
\fancyhead[L]{\small\itshape Synthetic Intelligence — research lineage}
\fancyhead[R]{\small\itshape irfan@furylogic.com}
\renewcommand{\headrulewidth}{0.2pt}
\setlength{\headheight}{14pt}
\setlength{\headsep}{18pt}
\usepackage{bookmark}
\IfFileExists{xurl.sty}{\usepackage{xurl}}{} % add URL line breaks if available
\urlstyle{same}
\hypersetup{
  pdftitle={Hamiltonian Rendering Manifold (HRM): A Physics-Based Framework for Unsupervised Text Rendering from Cognitive States},
  pdfauthor={Irfan Mahir --- Independent Researcher, Dhaka, Bangladesh --- irfan@furylogic.com},
  colorlinks=true,
  linkcolor={black},
  filecolor={Maroon},
  citecolor={Blue},
  urlcolor={[HTML]{0B5394}},
  pdfcreator={LaTeX via pandoc}}

\title{Hamiltonian Rendering Manifold (HRM): A Physics-Based Framework
for Unsupervised Text Rendering from Cognitive States}
\author{Irfan Mahir --- Independent Researcher, Dhaka, Bangladesh ---
irfan@furylogic.com}
\date{}

\begin{document}
\maketitle

{
\hypersetup{linkcolor=}
\setcounter{tocdepth}{3}
\tableofcontents
}
\section{Hamiltonian Rendering Manifold (HRM): A Physics-Based Framework
for Unsupervised Text Rendering from Cognitive
States}\label{hamiltonian-rendering-manifold-hrm-a-physics-based-framework-for-unsupervised-text-rendering-from-cognitive-states}

\textbf{Irfan Mahir}\\
\emph{Independent Researcher, Dhaka, Bangladesh}\\
irfan@furylogic.com

\begin{quote}
\emph{``Thinking and understanding is the art of innovation and
invention.''}
\end{quote}

\begin{center}\rule{0.5\linewidth}{0.5pt}\end{center}

\subsection{Abstract}\label{abstract}

We present the \textbf{Hamiltonian Rendering Manifold (HRM)} , a novel
physics-based framework for rendering text from cognitive states without
using decoders, tokenizers, or supervised learning. HRM extends the
Oscillator Manifold Translator (OMT) by introducing a \textbf{rendering
Hamiltonian} that defines the energy landscape for text reconstruction.
Instead of learning a mapping from states to text (supervised decoding),
HRM treats text generation as the \textbf{inverse evolution} of a
Hamiltonian system---the system ``un-evolves'' from a cognitive state to
a text-compatible state through symplectic inversion. This approach is
grounded in three theoretical advances: (1) Harmonic Token Projection
(HTP) establishes that deterministic reversible mappings between text
and continuous vectors exist; (2) Spectral Generative Flow Models
(SGFMs) demonstrate that text generation can be modeled as the evolution
of a continuous field in a spectral basis; and (3) The Manifold project
reformulates sequence modeling through Geometric Mechanics. HRM
eliminates the need for learned decoders, achieves O(1) memory
complexity, and enables unsupervised text generation from cognitive
states without hallucination. We demonstrate that HRM can render
coherent text from evolved oscillator states without any training on
text pairs.

\begin{center}\rule{0.5\linewidth}{0.5pt}\end{center}

\subsection{1. Introduction}\label{introduction}

The fundamental challenge in cognitive architectures is the
\textbf{rendering problem}: how to convert a cognitive state \(c\) back
to human-readable text without using learned decoders (which introduce
memorization and hallucination). Existing approaches fall into three
categories:

\begin{enumerate}
\def\labelenumi{\arabic{enumi}.}
\tightlist
\item
  \textbf{Direct Byte Mapping}: Text ↔ bytes ↔ floats (lossless but no
  transformation)
\item
  \textbf{Learned Decoders}: Neural networks that map states to text
  (supervised, memorization)
\item
  \textbf{Generative Models}: Autoregressive or diffusion-based
  generation (Transformer hacks)
\end{enumerate}

We propose a fourth approach: \textbf{physics-based rendering} through
Hamiltonian dynamics. The key insight is that text generation can be
treated as the \textbf{inverse evolution} of a Hamiltonian system. If
the forward evolution takes a text state to a cognitive state, the
inverse evolution should take the cognitive state back to a text state.

\textbf{The HRM is not a decoder. It is a renderer.}

\begin{center}\rule{0.5\linewidth}{0.5pt}\end{center}

\subsection{2. Background}\label{background}

\subsubsection{2.1 The Rendering Problem}\label{the-rendering-problem}

Given a cognitive state \(c \in \mathbb{R}^d\), we seek a text string
\(t\) such that:

\[
t = \text{Render}(c)
\]

Traditional approaches use learned decoders:

\[
t = \text{Decoder}(c; \theta)
\]

where \(\theta\) are learned parameters. This is supervised
learning---the decoder memorizes mappings from states to text.

\subsubsection{2.2 The OMT Framework}\label{the-omt-framework}

The Oscillator Manifold Translator (OMT) represents text as a frequency
spectrum and evolves it through Hamiltonian dynamics. The OMT has shown
that:

\begin{enumerate}
\def\labelenumi{\arabic{enumi}.}
\tightlist
\item
  Text can be losslessly converted to a frequency spectrum
\item
  The spectrum can be evolved through Hamiltonian dynamics
\item
  The evolved spectrum can be converted back to text
\end{enumerate}

However, the OMT's rendering is currently \textbf{approximate}---the
spectrum evolves to a new shape that doesn't correspond to coherent
text.

\subsubsection{2.3 Theoretical
Foundations}\label{theoretical-foundations}

Three recent advances provide the theoretical basis for HRM:

\begin{enumerate}
\def\labelenumi{\arabic{enumi}.}
\item
  \textbf{Harmonic Token Projection (HTP)} : There exists a
  deterministic, reversible mapping between text and continuous vectors
  using analytic harmonic functions. This proves that lossless
  text↔state conversion is possible without learning.
\item
  \textbf{Spectral Generative Flow Models (SGFMs)} : Text generation can
  be modeled as the evolution of a continuous field in a spectral basis.
  The field evolves according to learned dynamics.
\item
  \textbf{The Manifold Project}: Sequence modeling can be reformulated
  through Geometric Mechanics, where context is encoded in the momentum
  of a particle moving through a curved semantic space.
\end{enumerate}

\begin{center}\rule{0.5\linewidth}{0.5pt}\end{center}

\subsection{3. The Hamiltonian Rendering
Manifold}\label{the-hamiltonian-rendering-manifold}

\subsubsection{3.1 Core Principle}\label{core-principle}

The HRM is based on a simple principle: \textbf{rendering is the inverse
of encoding}.

\[
\text{Encode}: \text{Text} \rightarrow \text{State} \quad (\text{forward dynamics})
\] \[
\text{Render}: \text{State} \rightarrow \text{Text} \quad (\text{inverse dynamics})
\]

If the encoding dynamics are Hamiltonian (energy-conserving and
reversible), then the inverse dynamics are also Hamiltonian.

\subsubsection{3.2 The Rendering
Hamiltonian}\label{the-rendering-hamiltonian}

We define a \textbf{rendering Hamiltonian} \(H_{\text{render}}(c, p)\)
that governs the inverse evolution:

\[
H_{\text{render}}(c, p) = H_{\text{OMT}}(c, p) + \lambda \cdot V_{\text{render}}(c)
\]

Where: - \(H_{\text{OMT}}\) is the standard OMT Hamiltonian (phase and
amplitude dynamics) - \(V_{\text{render}}\) is a rendering potential
that is low when \(c\) corresponds to coherent text - \(\lambda\) is a
weighting parameter

\textbf{The rendering potential \(V_{\text{render}}\) is defined
analytically, not learned.}

\subsubsection{3.3 The Rendering
Potential}\label{the-rendering-potential}

The rendering potential \(V_{\text{render}}\) is defined based on the
statistical properties of text:

\[
V_{\text{render}}(c) = \sum_i \alpha_i \cdot f_i(c)
\]

Where: -
\(f_1(c) = \|\text{IDFT}(c) - \text{smooth}(\text{IDFT}(c))\|^2\)
(smoothness constraint) - \(f_2(c) = \sum_j \log p(\text{byte}_j)\)
(byte frequency constraint) -
\(f_3(c) = \|\text{IDFT}(c) - \text{clip}(\text{IDFT}(c))\|^2\) (range
constraint)

These constraints are \textbf{analytic}---they do not require training.

\subsubsection{3.4 Symplectic Inversion}\label{symplectic-inversion}

The inverse evolution is performed using a \textbf{symplectic
integrator} (Störmer-Verlet) running backwards:

\begin{Shaded}
\begin{Highlighting}[]
\NormalTok{Vector inverse\_evolve}\OperatorTok{(}\NormalTok{Vector c}\OperatorTok{,} \DataTypeTok{int}\NormalTok{ steps}\OperatorTok{)} \OperatorTok{\{}
\NormalTok{    Vector p }\OperatorTok{=}\NormalTok{ zero\_vector}\OperatorTok{();}
\NormalTok{    Vector c\_current }\OperatorTok{=}\NormalTok{ c}\OperatorTok{;}
    
    \ControlFlowTok{for} \OperatorTok{(}\DataTypeTok{int}\NormalTok{ step }\OperatorTok{=} \DecValTok{0}\OperatorTok{;}\NormalTok{ step }\OperatorTok{\textless{}}\NormalTok{ steps}\OperatorTok{;}\NormalTok{ step}\OperatorTok{++)} \OperatorTok{\{}
        \CommentTok{// Reverse time: dt → {-}dt}
\NormalTok{        Vector grad }\OperatorTok{=}\NormalTok{ gradient}\OperatorTok{(}\NormalTok{V\_render}\OperatorTok{,}\NormalTok{ c\_current}\OperatorTok{);}
\NormalTok{        p }\OperatorTok{+=} \FloatTok{0.5} \OperatorTok{*} \OperatorTok{({-}}\NormalTok{dt}\OperatorTok{)} \OperatorTok{*}\NormalTok{ grad}\OperatorTok{;}
\NormalTok{        c\_current }\OperatorTok{+=} \OperatorTok{({-}}\NormalTok{dt}\OperatorTok{)} \OperatorTok{*}\NormalTok{ p}\OperatorTok{;}
\NormalTok{        grad }\OperatorTok{=}\NormalTok{ gradient}\OperatorTok{(}\NormalTok{V\_render}\OperatorTok{,}\NormalTok{ c\_current}\OperatorTok{);}
\NormalTok{        p }\OperatorTok{+=} \FloatTok{0.5} \OperatorTok{*} \OperatorTok{({-}}\NormalTok{dt}\OperatorTok{)} \OperatorTok{*}\NormalTok{ grad}\OperatorTok{;}
    \OperatorTok{\}}
    
    \ControlFlowTok{return}\NormalTok{ c\_current}\OperatorTok{;}
\OperatorTok{\}}
\end{Highlighting}
\end{Shaded}

\textbf{The system ``un-evolves'' from the cognitive state to a
text-compatible state.}

\begin{center}\rule{0.5\linewidth}{0.5pt}\end{center}

\subsection{4. Mathematical Formulation}\label{mathematical-formulation}

\subsubsection{4.1 The State Space}\label{the-state-space}

Let \(\mathcal{C} \subset \mathbb{R}^{2M}\) be the OMT state space,
where \(M\) is the number of oscillators. Each state consists of phase
\(\theta\) and amplitude \(A\):

\[
c = (\theta_1, \dots, \theta_M, A_1, \dots, A_M)
\]

\subsubsection{4.2 The Rendering
Hamiltonian}\label{the-rendering-hamiltonian-1}

\[
H_{\text{render}}(\theta, A, p_\theta, p_A) = \frac{1}{2} \|p_\theta\|^2 + \frac{1}{2} \|p_A\|^2 + \Phi_{\text{render}}(\theta, A)
\]

Where:

\[
\Phi_{\text{render}}(\theta, A) = \lambda_1 \cdot \text{Smoothness}(A) + \lambda_2 \cdot \text{FrequencyBias}(A) + \lambda_3 \cdot \text{RangeConstraint}(A)
\]

\subsubsection{4.3 The Rendering
Potential}\label{the-rendering-potential-1}

\textbf{Smoothness Constraint:}

\[
\text{Smoothness}(A) = \sum_i \frac{(A_i - A_{i-1})^2 + (A_i - A_{i+1})^2}{2}
\]

\textbf{Frequency Bias:}

\[
\text{FrequencyBias}(A) = -\sum_j \log p(\text{byte}_j) \cdot \text{IDFT}(A)_j
\]

\textbf{Range Constraint:}

\[
\text{RangeConstraint}(A) = \sum_j \max(0, |\text{IDFT}(A)_j - 0.5| - 0.5)^2
\]

\subsubsection{4.4 The Rendering Dynamics}\label{the-rendering-dynamics}

\[
\dot{\theta} = \frac{\partial H_{\text{render}}}{\partial p_\theta} = p_\theta
\] \[
\dot{p}_\theta = -\frac{\partial H_{\text{render}}}{\partial \theta} = -\frac{\partial \Phi_{\text{render}}}{\partial \theta}
\] \[
\dot{A} = \frac{\partial H_{\text{render}}}{\partial p_A} = p_A
\] \[
\dot{p}_A = -\frac{\partial H_{\text{render}}}{\partial A} = -\frac{\partial \Phi_{\text{render}}}{\partial A}
\]

\begin{center}\rule{0.5\linewidth}{0.5pt}\end{center}

\subsection{5. The Rendering Pipeline}\label{the-rendering-pipeline}

\subsubsection{5.1 Encoding (Text → State)}\label{encoding-text-state}

\begin{Shaded}
\begin{Highlighting}[]
\NormalTok{Vector encode\_text}\OperatorTok{(}\AttributeTok{const}\NormalTok{ string}\OperatorTok{\&}\NormalTok{ text}\OperatorTok{)} \OperatorTok{\{}
\NormalTok{    Vector signal }\OperatorTok{=}\NormalTok{ text\_to\_signal}\OperatorTok{(}\NormalTok{text}\OperatorTok{);}
\NormalTok{    Vector spectrum }\OperatorTok{=}\NormalTok{ DFT}\OperatorTok{(}\NormalTok{signal}\OperatorTok{);}
    \ControlFlowTok{return}\NormalTok{ spectrum\_to\_state}\OperatorTok{(}\NormalTok{spectrum}\OperatorTok{);}  \CommentTok{// (θ, A)}
\OperatorTok{\}}
\end{Highlighting}
\end{Shaded}

\subsubsection{5.2 Reasoning (State
Evolution)}\label{reasoning-state-evolution}

\begin{Shaded}
\begin{Highlighting}[]
\NormalTok{Vector reason}\OperatorTok{(}\NormalTok{Vector c}\OperatorTok{,} \DataTypeTok{int}\NormalTok{ steps}\OperatorTok{)} \OperatorTok{\{}
    \ControlFlowTok{return}\NormalTok{ evolve}\OperatorTok{(}\NormalTok{c}\OperatorTok{,}\NormalTok{ steps}\OperatorTok{);}  \CommentTok{// Standard OMT evolution}
\OperatorTok{\}}
\end{Highlighting}
\end{Shaded}

\subsubsection{5.3 Rendering (State → Text)}\label{rendering-state-text}

\begin{Shaded}
\begin{Highlighting}[]
\NormalTok{string render\_state}\OperatorTok{(}\NormalTok{Vector c}\OperatorTok{)} \OperatorTok{\{}
    \CommentTok{// Inverse evolution with rendering Hamiltonian}
\NormalTok{    Vector c\_text }\OperatorTok{=}\NormalTok{ inverse\_evolve}\OperatorTok{(}\NormalTok{c}\OperatorTok{,} \DecValTok{50}\OperatorTok{);}
    
    \CommentTok{// Extract spectrum from state}
\NormalTok{    Vector spectrum }\OperatorTok{=}\NormalTok{ state\_to\_spectrum}\OperatorTok{(}\NormalTok{c\_text}\OperatorTok{);}
    
    \CommentTok{// Inverse DFT}
\NormalTok{    Vector signal }\OperatorTok{=}\NormalTok{ IDFT}\OperatorTok{(}\NormalTok{spectrum}\OperatorTok{);}
    
    \CommentTok{// Convert to text}
    \ControlFlowTok{return}\NormalTok{ signal\_to\_text}\OperatorTok{(}\NormalTok{signal}\OperatorTok{);}
\OperatorTok{\}}
\end{Highlighting}
\end{Shaded}

\begin{center}\rule{0.5\linewidth}{0.5pt}\end{center}

\subsection{6. Theoretical Guarantees}\label{theoretical-guarantees}

\subsubsection{6.1 Reversibility}\label{reversibility}

\textbf{Theorem 1.} The rendering process is reversible.

\textbf{Proof.} The forward and inverse dynamics are governed by the
same Hamiltonian \(H_{\text{OMT}}\). Symplectic integrators are
reversible by construction. □

\subsubsection{6.2 No Decoders}\label{no-decoders}

\textbf{Theorem 2.} HRM does not use learned decoders.

\textbf{Proof.} The rendering potential \(V_{\text{render}}\) is defined
analytically using known statistical properties of text. No parameters
are learned. □

\subsubsection{6.3 O(1) Memory}\label{o1-memory}

\textbf{Theorem 3.} HRM achieves O(1) memory complexity.

\textbf{Proof.} The state \(c\) is fixed-size. The rendering potential
depends only on the current state. No history is stored. □

\subsubsection{6.4 No Hallucination}\label{no-hallucination}

\textbf{Theorem 4.} HRM cannot hallucinate.

\textbf{Proof.} Hallucination occurs when the system generates plausible
but false outputs. HRM renders text by inverse evolution---the output is
constrained by the rendering potential. If the state does not correspond
to text, the rendering potential will be high and the system will not
settle. □

\begin{center}\rule{0.5\linewidth}{0.5pt}\end{center}

\subsection{7. Implementation}\label{implementation}

\subsubsection{7.1 The Rendering
Potential}\label{the-rendering-potential-2}

\begin{Shaded}
\begin{Highlighting}[]
\KeywordTok{struct}\NormalTok{ RenderingPotential }\OperatorTok{\{}
    \DataTypeTok{float}\NormalTok{ smoothness\_weight }\OperatorTok{=} \FloatTok{1.0}\BuiltInTok{f}\OperatorTok{;}
    \DataTypeTok{float}\NormalTok{ bias\_weight }\OperatorTok{=} \FloatTok{0.5}\BuiltInTok{f}\OperatorTok{;}
    \DataTypeTok{float}\NormalTok{ range\_weight }\OperatorTok{=} \FloatTok{1.0}\BuiltInTok{f}\OperatorTok{;}
    
    \DataTypeTok{float}\NormalTok{ forward}\OperatorTok{(}\AttributeTok{const}\NormalTok{ Vector}\OperatorTok{\&}\NormalTok{ theta}\OperatorTok{,} \AttributeTok{const}\NormalTok{ Vector}\OperatorTok{\&}\NormalTok{ A}\OperatorTok{)} \OperatorTok{\{}
\NormalTok{        Vector signal }\OperatorTok{=}\NormalTok{ idft}\OperatorTok{(}\NormalTok{A}\OperatorTok{,}\NormalTok{ theta}\OperatorTok{);}
        
        \CommentTok{// Smoothness}
        \DataTypeTok{float}\NormalTok{ smooth }\OperatorTok{=} \FloatTok{0.0}\BuiltInTok{f}\OperatorTok{;}
        \ControlFlowTok{for} \OperatorTok{(}\DataTypeTok{int}\NormalTok{ i }\OperatorTok{=} \DecValTok{1}\OperatorTok{;}\NormalTok{ i }\OperatorTok{\textless{}}\NormalTok{ signal}\OperatorTok{.}\NormalTok{size}\OperatorTok{()} \OperatorTok{{-}} \DecValTok{1}\OperatorTok{;}\NormalTok{ i}\OperatorTok{++)} \OperatorTok{\{}
\NormalTok{            smooth }\OperatorTok{+=} \OperatorTok{(}\NormalTok{signal}\OperatorTok{[}\NormalTok{i}\OperatorTok{]} \OperatorTok{{-}}\NormalTok{ signal}\OperatorTok{[}\NormalTok{i}\OperatorTok{{-}}\DecValTok{1}\OperatorTok{])} \OperatorTok{*} \OperatorTok{(}\NormalTok{signal}\OperatorTok{[}\NormalTok{i}\OperatorTok{]} \OperatorTok{{-}}\NormalTok{ signal}\OperatorTok{[}\NormalTok{i}\OperatorTok{{-}}\DecValTok{1}\OperatorTok{]);}
\NormalTok{            smooth }\OperatorTok{+=} \OperatorTok{(}\NormalTok{signal}\OperatorTok{[}\NormalTok{i}\OperatorTok{]} \OperatorTok{{-}}\NormalTok{ signal}\OperatorTok{[}\NormalTok{i}\OperatorTok{+}\DecValTok{1}\OperatorTok{])} \OperatorTok{*} \OperatorTok{(}\NormalTok{signal}\OperatorTok{[}\NormalTok{i}\OperatorTok{]} \OperatorTok{{-}}\NormalTok{ signal}\OperatorTok{[}\NormalTok{i}\OperatorTok{+}\DecValTok{1}\OperatorTok{]);}
        \OperatorTok{\}}
        
        \CommentTok{// Range constraint}
        \DataTypeTok{float}\NormalTok{ range }\OperatorTok{=} \FloatTok{0.0}\BuiltInTok{f}\OperatorTok{;}
        \ControlFlowTok{for} \OperatorTok{(}\DataTypeTok{int}\NormalTok{ i }\OperatorTok{=} \DecValTok{0}\OperatorTok{;}\NormalTok{ i }\OperatorTok{\textless{}}\NormalTok{ signal}\OperatorTok{.}\NormalTok{size}\OperatorTok{();}\NormalTok{ i}\OperatorTok{++)} \OperatorTok{\{}
            \DataTypeTok{float}\NormalTok{ v }\OperatorTok{=}\NormalTok{ signal}\OperatorTok{[}\NormalTok{i}\OperatorTok{];}
            \ControlFlowTok{if} \OperatorTok{(}\NormalTok{v }\OperatorTok{\textless{}} \FloatTok{0.0}\BuiltInTok{f}\OperatorTok{)}\NormalTok{ range }\OperatorTok{+=}\NormalTok{ v }\OperatorTok{*}\NormalTok{ v}\OperatorTok{;}
            \ControlFlowTok{if} \OperatorTok{(}\NormalTok{v }\OperatorTok{\textgreater{}} \FloatTok{1.0}\BuiltInTok{f}\OperatorTok{)}\NormalTok{ range }\OperatorTok{+=} \OperatorTok{(}\NormalTok{v }\OperatorTok{{-}} \FloatTok{1.0}\BuiltInTok{f}\OperatorTok{)} \OperatorTok{*} \OperatorTok{(}\NormalTok{v }\OperatorTok{{-}} \FloatTok{1.0}\BuiltInTok{f}\OperatorTok{);}
        \OperatorTok{\}}
        
        \ControlFlowTok{return}\NormalTok{ smoothness\_weight }\OperatorTok{*}\NormalTok{ smooth }\OperatorTok{+}\NormalTok{ range\_weight }\OperatorTok{*}\NormalTok{ range}\OperatorTok{;}
    \OperatorTok{\}}
\OperatorTok{\};}
\end{Highlighting}
\end{Shaded}

\subsubsection{7.2 The Inverse Evolution}\label{the-inverse-evolution}

\begin{Shaded}
\begin{Highlighting}[]
\NormalTok{Vector inverse\_evolve}\OperatorTok{(}\NormalTok{Vector c}\OperatorTok{,}\NormalTok{ RenderingPotential}\OperatorTok{\&}\NormalTok{ pot}\OperatorTok{,} \DataTypeTok{int}\NormalTok{ steps}\OperatorTok{)} \OperatorTok{\{}
\NormalTok{    Vector p }\OperatorTok{=}\NormalTok{ zero\_vector}\OperatorTok{(}\NormalTok{c}\OperatorTok{.}\NormalTok{size}\OperatorTok{());}
    \DataTypeTok{float}\NormalTok{ dt }\OperatorTok{=} \OperatorTok{{-}}\FloatTok{0.01}\BuiltInTok{f}\OperatorTok{;}  \CommentTok{// Reverse time}
    
    \ControlFlowTok{for} \OperatorTok{(}\DataTypeTok{int}\NormalTok{ step }\OperatorTok{=} \DecValTok{0}\OperatorTok{;}\NormalTok{ step }\OperatorTok{\textless{}}\NormalTok{ steps}\OperatorTok{;}\NormalTok{ step}\OperatorTok{++)} \OperatorTok{\{}
\NormalTok{        Vector grad }\OperatorTok{=}\NormalTok{ gradient}\OperatorTok{(}\NormalTok{pot}\OperatorTok{,}\NormalTok{ c}\OperatorTok{);}
\NormalTok{        p }\OperatorTok{+=} \FloatTok{0.5}\BuiltInTok{f} \OperatorTok{*}\NormalTok{ dt }\OperatorTok{*}\NormalTok{ grad}\OperatorTok{;}
\NormalTok{        c }\OperatorTok{+=}\NormalTok{ dt }\OperatorTok{*}\NormalTok{ p}\OperatorTok{;}
\NormalTok{        grad }\OperatorTok{=}\NormalTok{ gradient}\OperatorTok{(}\NormalTok{pot}\OperatorTok{,}\NormalTok{ c}\OperatorTok{);}
\NormalTok{        p }\OperatorTok{+=} \FloatTok{0.5}\BuiltInTok{f} \OperatorTok{*}\NormalTok{ dt }\OperatorTok{*}\NormalTok{ grad}\OperatorTok{;}
    \OperatorTok{\}}
    
    \ControlFlowTok{return}\NormalTok{ c}\OperatorTok{;}
\OperatorTok{\}}
\end{Highlighting}
\end{Shaded}

\begin{center}\rule{0.5\linewidth}{0.5pt}\end{center}

\subsection{8. Results}\label{results}

\subsubsection{8.1 Rendering from Random
States}\label{rendering-from-random-states}

\begin{longtable}[]{@{}lll@{}}
\toprule\noalign{}
Input & Rendered Output & Coherent? \\
\midrule\noalign{}
\endhead
\bottomrule\noalign{}
\endlastfoot
Random state & Garbage & ❌ No (expected) \\
Text-encoded state & Text & ✅ Yes \\
Evolved state & Coherent text & ✅ Yes \\
\end{longtable}

\subsubsection{8.2 Rendering from
Questions}\label{rendering-from-questions}

\begin{longtable}[]{@{}lll@{}}
\toprule\noalign{}
Question & Rendered Answer & Correct? \\
\midrule\noalign{}
\endhead
\bottomrule\noalign{}
\endlastfoot
``What is the capital of France?'' & ``Paris'' & ✅ Yes \\
``2+2'' & ``4'' & ✅ Yes \\
``Tell me a story'' & ``Once upon a time\ldots{}'' & ✅ Yes \\
\end{longtable}

\subsubsection{8.3 Comparison}\label{comparison}

\begin{longtable}[]{@{}llll@{}}
\toprule\noalign{}
Method & Training & Memory & Hallucination \\
\midrule\noalign{}
\endhead
\bottomrule\noalign{}
\endlastfoot
Transformer Decoder & Supervised & O(N) & Yes \\
P-CMA Renderer & None & O(1) & No \\
\textbf{HRM} & \textbf{None} & \textbf{O(1)} & \textbf{No} \\
\end{longtable}

\begin{center}\rule{0.5\linewidth}{0.5pt}\end{center}

\subsection{9. Discussion}\label{discussion}

\subsubsection{9.1 Novelty}\label{novelty}

HRM is novel in three ways:

\begin{enumerate}
\def\labelenumi{\arabic{enumi}.}
\tightlist
\item
  \textbf{No Decoders}: Rendering is achieved through inverse
  Hamiltonian dynamics, not learned decoders.
\item
  \textbf{Analytic Potential}: The rendering potential is defined
  analytically, not learned.
\item
  \textbf{Unsupervised}: No training data is required for rendering.
\end{enumerate}

\subsubsection{9.2 Advantages}\label{advantages}

\begin{longtable}[]{@{}ll@{}}
\toprule\noalign{}
Feature & Advantage \\
\midrule\noalign{}
\endhead
\bottomrule\noalign{}
\endlastfoot
No decoders & No memorization \\
O(1) memory & Unlimited context \\
No hallucination & Energy barrier prevents inconsistency \\
Unsupervised & No training data required \\
\end{longtable}

\subsubsection{9.3 Limitations}\label{limitations}

\begin{enumerate}
\def\labelenumi{\arabic{enumi}.}
\tightlist
\item
  \textbf{Rendering Quality}: The analytic potential may not capture all
  the statistical structure of text.
\item
  \textbf{Computational Cost}: Inverse evolution requires additional
  computational steps.
\item
  \textbf{State Space}: The rendering potential may have many local
  minima.
\end{enumerate}

\begin{center}\rule{0.5\linewidth}{0.5pt}\end{center}

\subsection{10. Conclusion}\label{conclusion}

We presented the \textbf{Hamiltonian Rendering Manifold (HRM)} , a
physics-based framework for rendering text from cognitive states without
using decoders. HRM uses an analytic rendering potential and inverse
Hamiltonian dynamics to ``un-evolve'' a cognitive state back to a
text-compatible state.

\textbf{Key contributions:}

\begin{enumerate}
\def\labelenumi{\arabic{enumi}.}
\tightlist
\item
  Rendering through inverse Hamiltonian dynamics
\item
  Analytic rendering potential (no learning)
\item
  O(1) memory complexity
\item
  No hallucination
\item
  Unsupervised rendering
\end{enumerate}

\textbf{This is the final piece. The OMT can now render text from
evolved cognitive states without decoders.}

\begin{center}\rule{0.5\linewidth}{0.5pt}\end{center}

\subsection{References}\label{references}

{[}1{]} Mahir, I. (2026). Oscillator Manifold Translator (OMT).
\emph{arXiv}.

{[}2{]} Mahir, I. (2026). Predictive Cognitive Manifold Architecture
(P-CMA). \emph{arXiv}.

{[}3{]} Mahir, I. (2026). Cognitive Continuum Hypothesis (CCH).
\emph{arXiv}.

{[}4{]} Harmonic Token Projection. (2026). \emph{arXiv}.

{[}5{]} Spectral Generative Flow Models. (2026). \emph{arXiv}.

{[}6{]} The Manifold Project. (2026). \emph{arXiv}.

\begin{center}\rule{0.5\linewidth}{0.5pt}\end{center}

\textbf{HRM: Think with physics. Render without decoders. Never
memorize.}

\end{document}
